\documentclass[11pt,a4paper]{article}
\usepackage[utf8]{inputenc}
\usepackage[T1]{fontenc}
\usepackage{lmodern}
\usepackage{xcolor}
\renewcommand{\contentsname}{Spis treści}
\renewcommand{\tablename}{Tabela}
\hyphenpenalty=3000 \tolerance=2000 \emergencystretch=2em
\usepackage[margin=2.2cm]{geometry}
\usepackage{booktabs}
\usepackage{tabularx}
\usepackage{longtable}
\usepackage{enumitem}
\usepackage[colorlinks=true,linkcolor=blue!50!black,urlcolor=blue!50!black]{hyperref}
\usepackage{titlesec}
\usepackage{parskip}
\titleformat{\section}{\large\bfseries}{\thesection.}{0.6em}{}
\titleformat{\subsection}{\normalsize\bfseries}{\thesubsection.}{0.6em}{}
\setlist{nosep,leftmargin=1.4em}
\newcommand{\ohm}{$\Omega$}

\title{\textbf{AMPERE POINT --- Wielofunkcyjne urządzenie pomiarowe (custom device)}\\[4pt]
\large Koncepcja budowy --- wersja 5\\
budżetowy \textbf{elektryczny} magazyn energii (używane prostowniki + bateria składana samodzielnie)}
\author{Opracowanie wewnętrzne AMPERE POINT}
\date{8 lipca 2026}

\begin{document}
\maketitle
\vspace{-1.5em}

\noindent\textbf{Status dokumentu:} rewizja różnicowa. Decyzja zamawiającego: magazyn ma być
\textbf{elektryczny} (uniwersalny --- zasila cokolwiek i kiedykolwiek, bez wymogu bieżącego zużycia,
jak w termicznym v4), a niski koszt osiągamy \textbf{własną pracą montażową} zamiast zakupem
gotowych klocków. Wersja 5 zastępuje v4; architektura elektryczna pozostaje ta z v2/v3
(PFC $\to$ szyna DC 48~V $\to$ bateria $\to$ falownik), więc \textbf{schemat elektryczny v4 (EPLAN)
pozostaje aktualny} po podmianie typów aparatów (rozdz.~\ref{sec:dok}). Moduł A, sekwencja,
łańcuch bezpieczeństwa i decyzje Z1--Z6 --- bez zmian. Nowe decyzje: Z16--Z19.

\tableofcontents

\section{Zasada cięcia kosztów: ,,tanio przez pracę''}
Trzy najdroższe pozycje rekuperacji z v3 zastępujemy odpowiednikami z rynku wtórnego i montażem
własnym:
\begin{enumerate}
\item \textbf{Prostowniki PFC}: zamiast 6$\times$ Huawei R4850G2 ($\sim$2100--3000~zł) ---
6$\times$ \textbf{Eltek Flatpack2 48/2000 HE} z demontaży telekomunikacyjnych
(\textbf{$\sim$150--300~zł/szt.}), sprawność 96{,}5\%, PF$>$0,99, \textbf{sterowanie prądem po CAN}
(protokół dobrze opisany przez społeczność, lepiej niż R4850). 6$\times$2~kW $=$ 12~kW.
\item \textbf{Bateria}: zamiast gotowego racka LiFePO4 (4000--6500~zł) --- \textbf{pakiet 48~V
składany samodzielnie} z ogniw LFP klasy~B / z odzysku albo z używanych modułów EV
(rozdz.~\ref{sec:bat}); oszczędność 30--50\% za cenę 1--2 dni pracy.
\item \textbf{Falownik}: warianty W2 (używany UPS on-line z szyną 48~V, 500--1500~zł) lub W3
(własny LF wg gotowego przepisu z koncepcji v3, rozdz.~6; 1350--2490~zł) --- praca wyspowa na
wydzielone obwody; W1 (hybrydowy, ESS) tylko jako przyszły upgrade.
\end{enumerate}
Dodatkowe cięcia: szyny DC z płaskownika Cu (zamiast systemowych), montaż prostowników na
własnym wsporniku (używana półka Flatpack tylko jeśli trafi się tanio), \textbf{rekonfiguracja
faz na sztywnym polu przełączeń (ręczne mostki przed testem 1F 32~A) zamiast stycznika -K14} ---
sekwencer i tak weryfikuje konfigurację pomiarem (-U2) przed zamknięciem -K1.

\section{Architektura (jak v2/v3, z podmianami)}
\begin{center}\small
DUT $\to$ moduł A (tor mocy, -B1/-T/-K1) $\to$ CEE $\to$ -K11..-K13 $\to$
\textbf{6$\times$ Flatpack2 48/2000 HE} (CAN) $\to$ szyna DC 48~V (-F13, -KB, -Q1) $\to$
\textbf{-GB1: pakiet 48~V DIY} $\to$ \textbf{falownik W2/W3} $\to$ wydzielone obwody warsztatu.
\end{center}
Przydział mocy: 3F 16~A $=$ 2 szt./fazę (4~kW $>$ 3,7~kW) (spełnione); 1F 32~A $=$ 4 szt. na L1 (8~kW $>$
7,4~kW) --- dwie sztuki przepinane mostkami na polu przełączeń. Pobór sinusoidalny, PF$>$0,99,
rampy i kroki 25/50/75/100\% deklaracji CP zadawane po CAN (ESP32-S3/TWAI, jak w v3).
\textbf{Bez zrzutu cieplnego}: zamiast grzałki CWU --- \textbf{bramka programowa SOC}
(rozdz.~\ref{sec:soc}); jedyna ,,grzałka'' w systemie odpada zupełnie.

\section{Bateria -GB1: warianty ogniw i przepis montażu}
\label{sec:bat}
\subsection{Wymiar energetyczny}
Soak 30~min przy 11~kW $=$ 5,5~kWh; przy oknie roboczym SOC 40$\to$90\% potrzeba $\sim$11~kWh
pojemności na pełny soak ,,z zapasem''. Strategia budżetowa: \textbf{start od jednej gałęzi 16S
(5,4~kWh)} --- wystarcza na soak 1F (3,7~kWh) i skrócone/łagodniejsze soaki 3F, zwłaszcza że
falownik może \emph{równolegle} rozładowywać pakiet na odbiory --- a \textbf{rozbudowa do 2P
(10,8~kWh) to tylko dokupienie 16 ogniw} i dołożenie drugiej gałęzi (architektura bez zmian).

\subsection{Warianty ogniw (decyzja Z17)}
\begin{longtable}{@{}p{0.20\textwidth}p{0.26\textwidth}p{0.22\textwidth}p{0.24\textwidth}@{}}
\toprule
 & \textbf{B-a: moduły EV z odzysku} & \textbf{B-b: ogniwa LFP klasy B (rekomendacja)} & \textbf{B-c: gotowy rack 48 V}\\
\midrule
przykład & moduły NMC z demontażu (i3, Leaf itp.) & 16$\times$ EVE LF105 / CATL 100--105 Ah & rack 10 kWh z BMS\\
koszt / $\sim$10 kWh & 2500--4000 zł & 4500--6400 zł (2P); 1P 5,4 kWh: 2240--3200 zł & 4000--6500 zł\\
chemia / bezpieczeństwo & NMC --- wyższe ryzyko termiczne: stalowy przedział, czujka dymu, ładowanie do $\sim$80\% & LFP --- stabilna termicznie, właściwa do warsztatu & LFP\\
praca własna & duża (adaptacja modułów, BMS) & średnia (montaż pakietu wg 3.3) & minimalna\\
\bottomrule
\end{longtable}

\subsection{Przepis montażu pakietu 16S LFP (B-b) krok po kroku}
\begin{enumerate}
\item \textbf{Zakup i selekcja.} 16 (docelowo 32) ogniw 100--105 Ah klasy B z pomiarami sprzedawcy
(pojemność, IR); po odbiorze zmierz napięcia (rozrzut $\le$30 mV) i rezystancję wewnętrzną
(UT890C + metoda przyrostu napięcia albo tani tester YR1035); paruj ogniwa pojemnościami.
\item \textbf{Top-balance.} Połącz wszystkie ogniwa \emph{równolegle} i ładuj do 3,55--3,60 V
źródłem CV --- \textbf{trik: jeden Flatpack2 ustawiony po CAN na 3,55 V $\times$ 16 $=$ 56,8 V
posłuży później; do balansu równoległego użyj zasilacza lab./ładowarki LFP 3,6 V} (albo balansuj
gałąź po złożeniu funkcją BMS --- dłużej). Czas: 12--48 h; to najnudniejszy i najważniejszy krok.
\item \textbf{Kompresja.} Płyty czołowe (sklejka 18 mm/stal) + 4 pręty M8: docisk zgodny z notą
ogniw ($\sim$300 kgf dla LF105) --- wydłuża żywotność i usztywnia pakiet.
\item \textbf{Połączenia.} Szyny fabryczne ogniw lub płaskownik Cu; momenty wg noty (typ. 4--6 Nm
na M6); po skręceniu pomiar spadków na łączach przy próbnym prądzie (UT210E + UT890C).
\item \textbf{BMS.} Smart BMS 16S 150--200 A z balansem aktywnym 1--2 A (klasy JK/JBD, 350--600 zł):
przewody balansera podłączaj \textbf{od minusa, kolejno}, wtyk na końcu; konfiguracja: 2,80/3,55 V
na ogniwo, prąd 150 A, ładowanie 0--45 $^\circ$C, 2$\times$NTC na ogniwach skrajnych;
styk/wyjście BMS $\to$ kontaktor -KB (odcięcie niezależne od MCU, jak w v3).
\item \textbf{Zabezpieczenia gałęzi.} Bezpiecznik mega 150--200 A przy biegunie $+$, przewody
50 mm$^2$ zakończone prasowanymi końcówkami M8; rozłącznik -Q1 w torze.
\item \textbf{Pierwsze ładowanie nadzorowane.} Z jednego Flatpacka (prąd ograniczony do 10--20 A),
obserwacja napięć ogniw w aplikacji BMS; test zadziałania progów (podnieś próg, sprawdź odcięcie).
\item \textbf{Pomiar pojemności kontrolny.} Rozładowanie przez falownik na znane obciążenie
(np. 1--2 kW) do progu BMS z rejestracją w -A1 --- wynik do dokumentacji pakietu.
\item \textbf{Zabudowa.} Dolny przedział szafy, przegroda od strefy AC, wentylacja grawitacyjna,
czujniki DS18B20 na ogniwach do 1-Wire (-A1); etykieta: pojemność, data, progi.
\end{enumerate}

\section{Zarządzanie energią bez zrzutu cieplnego}
\label{sec:soc}
W v5 nie ma grzałki zrzutowej --- rolę ,,zaworu'' pełni \textbf{bramka SOC} w sekwencerze:
przed soakiem 3F wymagane SOC $\le$40\% (1P: $\le$25\% lub soak skrócony; 2P: $\le$40\%),
przed soakiem 1F $\le$60\%; przy niespełnieniu HMI komunikuje ,,rozładuj magazyn'' (włącz odbiory
wyspy / poczekaj). Podczas testu falownik może pracować --- rozładowanie równoległe zmniejsza
przyrost SOC netto. Telemetria: SOC/prądy z BMS po RS485/CAN do -A1; niezależnie -U2 liczy energię
pobraną z DUT. Opcja za $\sim$100 zł (plan C, nieobowiązkowy): rezystor awaryjny 2 kW załączany
ręcznie do szybkiego zrzutu przed pilnym testem.

\section{Kosztorys modułu B v5 i porównanie wariantów}
\begin{longtable}{@{}p{0.58\textwidth}r@{}}
\toprule
\textbf{Pozycja} & \textbf{zł}\\
\midrule
6$\times$ Eltek Flatpack2 48/2000 HE (rynek wtórny) & 900--1800\\
wspornik/półka + złącza wyjściowe prostowników & 100--400\\
-K11..-K13, rozdział AC, pole przełączeń 1F (mostki) & 350--600\\
szyna DC (płaskownik Cu), -F13.1/2, -Q1, -KB, przewody 50 mm$^2$ & 600--1100\\
bateria B-b1: 16S1P LF105 kl. B + BMS + kompresja/szyny & 2740--4100\\
\quad (rozbudowa B-b2 do 2P 10,8 kWh: dodatkowo) & (+2240--3200)\\
falownik W2 (używany UPS on-line 48 V) \emph{lub} W3 (własny LF wg v3) & 500--2490\\
szafa, przegrody, wentylatory strat, drobnica & 400--800\\
\midrule
\textbf{Moduł B v5 razem (B-b1 + W2)} & \textbf{$\sim$5600--10300}\\
\textbf{Moduł B v5 razem (B-b2 + W2/W3)} & \textbf{$\sim$7800--13500}\\
\bottomrule
\end{longtable}

\begin{longtable}{@{}p{0.26\textwidth}rp{0.16\textwidth}p{0.34\textwidth}@{}}
\toprule
\textbf{Wariant modułu B} & \textbf{koszt [zł]} & \textbf{magazyn} & \textbf{uwagi}\\
\midrule
v1: grzałki + wentylacja & 2600--4600 & brak & energia tracona\\
v4: zasobnik wodny & 2150--4780 & ciepło 21--35 kWh & wymaga zużycia ciepła (odrzucone)\\
\textbf{v5: magazyn elektryczny DIY} & \textbf{5600--13500} & \textbf{prąd 5,4--10,8 kWh} & uniwersalny; praca własna 2--4 dni\\
v3: rekuperacja ,,z półki'' & 10850--18500 & prąd 10--14 kWh & najmniej pracy, najdrożej\\
\midrule
\textbf{Stacja razem (moduł A + B v5, B-b1+W2)} & \textbf{$\sim$9350--17350} & & v3 było $\sim$14200--27050\\
\bottomrule
\end{longtable}
Uczciwie: magazynu elektrycznego nie da się sprowadzić do cen v1/v4 --- bateria jest kosztem
nieredukowalnym; praca własna tnie wszystko wokół niej ($\sim$40--45\% względem v3).

\section{Bezpieczeństwo, ryzyka, etapy}
Strona DC jak w v3 (przekroje 50--70 mm$^2$, -F13, -Q1, -KB niezależny od MCU, momenty połączeń);
przedział bateryjny z przegrodą; przy wariancie B-a (NMC) dodatkowo stalowa obudowa i czujka dymu.
Nowe ryzyka: \textbf{R20} jakość ogniw used (selekcja/IR/parowanie; kupować z pomiarami),
\textbf{R21} NMC termicznie (rekomendacja LFP), \textbf{R22} protokół CAN Flatpack2
(dokumentacja społecznościowa --- weryfikacja w etapie 0; fallback R4850G2), \textbf{R23} czas
pracy własnej (2--4 dni pakiet + szafa), \textbf{R24} 1P za mały na pełny soak 3F (bramka SOC,
rozładowanie równoległe, rozbudowa 2P). Etapy: \textbf{0} --- 1$\times$Flatpack2 + sterowanie CAN
z ESP32 + stanowisko top-balance; \textbf{1} --- 3--4$\times$FP2 na L1 + pakiet 1P + W2: pełny
test 1F z magazynem; \textbf{2} --- 6$\times$FP2, pole przełączeń, rozbudowa 2P (3F/Q11);
\textbf{3} --- ewentualny upgrade falownika do W1 (ESS) --- bez zmian reszty.

\section{Wpływ na dokumentację}
\label{sec:dok}
\textbf{Schemat elektryczny v4 (EPLAN) pozostaje bazą} --- zmiany do naniesienia w v5 schematu
(po Z16): typy -PR1..-PR6 $=$ Flatpack2 48/2000 HE; -K14 $\to$ pole przełączeń ręcznych
(adnotacja + kontrola konfiguracji pomiarem -U2); usunięcie -K5/-E13 (zrzut CWU) --- zastępuje je
bramka SOC w firmware. Wizualizacja 3D v2 (szafa rekuperacji) pozostaje reprezentatywna dla v5.
Koncepcje v3/v4 zostają w folderze jako warianty odłożone.

\section{Punkty do zatwierdzenia}
\begin{enumerate}[label=\textbf{Z\arabic*.},start=16]
\item \textbf{Kierunek v5}: budżetowy magazyn elektryczny (używane FP2 + pakiet DIY + W2/W3),
zastępuje v4. \emph{Rekomendacja: tak.}
\item \textbf{Ogniwa}: B-b LFP klasy B (rekomendacja), start 1P 5,4 kWh $\to$ rozbudowa 2P;
alternatywy B-a (EV/NMC, taniej, ostrożnie) / B-c (rack, drożej).
\item \textbf{Falownik}: W2 --- używany UPS (rekomendacja: najmniej pracy za najmniejsze
pieniądze) vs W3 --- własny LF (przepis w v3).
\item \textbf{Budżet stacji} $\sim$9,4--17,4 tys. zł (B-b1+W2). \emph{Do akceptacji.}
\end{enumerate}

\section{Źródła (uzupełnienie)}
\begin{itemize}\small
\item Eltek Flatpack2 48/2000 HE --- karta produktu; protokół CAN: dokumentacja społecznościowa
(projekty openinverter/endless-sphere).
\item EVE LF105 / CATL --- noty ogniw LFP (kompresja, progi napięć, temperatury ładowania).
\item Smart BMS 16S klasy JK/JBD --- instrukcje konfiguracji, balans aktywny.
\item Pozostałe: jak koncepcja v3 (EG8010/EGS002, PN-EN 50549-1, IEC 62955, PN-HD 60364).
\end{itemize}

\end{document}
